\documentclass[10pt,a4paper,landscape]{article}

\usepackage[margin=1.25cm]{geometry}
\usepackage[T1]{fontenc}
\usepackage{lmodern}
\usepackage{array}
\usepackage{booktabs}
\usepackage{longtable}
\usepackage[table]{xcolor}
\usepackage[hidelinks]{hyperref}
\usepackage{microtype}

\setlength{\parindent}{0pt}
\setlength{\parskip}{0.45em}
\setlength{\tabcolsep}{4pt}
\renewcommand{\arraystretch}{1.18}
\newcommand{\code}[1]{\texttt{#1}}
\newcommand{\commit}[1]{\code{#1}}
\newcommand{\gainmeasured}{\textbf{Measured}}
\newcommand{\gainbounded}{\textbf{Bounded}}
\newcommand{\gainenabling}{\textbf{Enabling}}
\newcommand{\gainnoise}{\textbf{No clear gain}}
\newcommand{\stateactive}{\textbf{Active}}
\newcommand{\stateoptional}{\textbf{Optional}}
\newcommand{\statereverted}{\textbf{Reverted}}

\newcolumntype{I}{>{\raggedright\arraybackslash}p{2.3cm}}
\newcolumntype{O}{>{\raggedright\arraybackslash}p{4.1cm}}
\newcolumntype{R}{>{\raggedright\arraybackslash}p{4.4cm}}
\newcolumntype{G}{>{\raggedright\arraybackslash}p{4.6cm}}
\newcolumntype{K}{>{\raggedright\arraybackslash}p{5.2cm}}
\newcolumntype{S}{>{\raggedright\arraybackslash}p{2.0cm}}

\title{cMVBT Optimization Catalog\\\large Complete repository inventory, gains, reasons, and risks}
\author{Repository documentation}
\date{11 August 2026}

\begin{document}
\maketitle

\section*{Purpose and evidence rules}

This document is the compact index of optimizations found in the cMVBT repository and
its Git history. It covers the current tree, earlier optimizations that shaped it, and
important attempted optimizations that were rejected or reverted. Correctness fixes are
included when they remove retry, livelock, overflow, leak, or unbounded-growth behaviour;
those are operational optimizations even when no throughput number was recorded.

The gain column deliberately distinguishes three evidence levels:

\begin{itemize}
  \item \gainmeasured{}: a before/after benchmark or profiler measurement exists in the
        repository documentation. Numbers retain the original workload and caveats.
  \item \gainbounded{}: the implementation removes or bounds a countable cost (allocations,
        tuples, locks, retries, scans, or memory growth), but no trustworthy end-to-end
        percentage was recorded.
  \item \gainenabling{}: primarily a correctness or scalability fix. It enables completion,
        safe reclamation, or correct concurrency rather than claiming a numeric speedup.
\end{itemize}

Combined benchmark results must not be attributed to each component independently. For
example, the first WAL result table measures CRC, build, and channel-sharding changes as
one set. ``Unknown'' never means zero; it means the repository contains no controlled
measurement from which a number can responsibly be inferred.

Primary detailed sources are \code{docs/optimizations.tex},
\code{docs/oltp\_wal\_optimization.md}, \code{docs/range\_scan\_iteration.md},
\code{docs/range\_scan\_visibility\_check.md}, and
\code{docs/bigtree\_size\_benchmark.md}.

\section{Concurrency, synchronization, and memory ordering}

\begin{longtable}{IORGKS}
\toprule
\textbf{Commit / location} & \textbf{Optimization} & \textbf{Reason} &
\textbf{Gain / evidence} & \textbf{Potential risk or trade-off} & \textbf{State} \\
\midrule
\endhead
\commit{9abfb51}, current OLC paths & Optimistic lock coupling: lock-free readers and
version-validated traversal & Avoid reader latches and hand-over-hand lock traffic;
writers use a seqlock-like node word. & \gainenabling: unlimited concurrent readers; current
read path takes an Acquire version sample rather than a latch. Numeric historical gain
was not preserved. & Every field observed without a latch must be atomic or immutable;
missed revalidation becomes a use-after-reuse or torn-read bug. & \stateactive \\

\commit{95cc76d} & Remove the master/root guard from write traversal; use atomic state
checks & The global master guard serialized otherwise independent writes. & \gainbounded:
one global guard removed from each write traversal. & Root replacement must still return
a pinned, validated root; stale unsafe-state observations can split the wrong root. &
\stateactive \\

\commit{f85d6d1}, \commit{e5b07a0}, \commit{867f28e} & OLC borrow and traversal
micro-optimizations & Reduce wrapper work and repeated validation on point traversals. &
\gainbounded: fewer borrow/guard operations; no retained controlled percentage. & Extremely
hot and correctness-sensitive path; a removed validation may admit a retired block. &
\stateactive \\

\commit{48c4ded}, \commit{3c6da7b}, \commit{a518eed} & Remove standalone fences already
subsumed by atomic load/CAS ordering & Earlier defensive fences duplicated ordering in
the adjacent atomic operation. & \gainbounded: removes fence instructions from traversal and
page-length paths. & Weakening ordering without a complete happens-before argument can
make initialized page data invisible to readers. & \stateactive \\

\commit{a7bb29c} & Use Acquire/Release semantics that match OLC publication & Avoid
unnecessary CAS failures while establishing the exact reader/writer edge. & \gainbounded:
fewer spurious retries; no isolated benchmark retained. & Incorrect success/failure
ordering silently breaks the seqlock contract on weakly ordered CPUs. & \stateactive \\

\commit{2a2be93} & Page soft-commit counters use Relaxed stores & The later latch unlock
already performs the Release publication. & \gainbounded: one weaker atomic operation per
page commit. & Safe only while every path performs the releasing latch unlock after the
counter update. & \stateactive \\

\commit{40b00fd} & Make the node type tag non-atomic and remove the Drop fence & The tag
changes only under the exclusive node latch. & \gainbounded: removes atomic/fence work and
shrinks metadata. & Any future out-of-latch mutation would introduce a data race. &
\stateactive \\

\commit{45c7c21}/\commit{fd92f38} & Direct \code{try\_retire} instead of a conservative
child write-latch upgrade & A guard that is unconditionally consumed does not need a
reversible exclusive upgrade. & \gainbounded: one CAS/write-lock acquisition removed per
eligible split or merge child. & Unsound for any path that can back out after retirement;
such paths must explicitly unretire or retain the full upgrade. & \stateactive \\

\commit{7f2bd09} & Cache-line-align the hot node tag and use SmartCell for root liveness &
Avoid false sharing and a separate root mutex. & \gainbounded: isolates the frequently polled
tag; no controlled throughput number. & Padding increases object size; alignment only
helps if allocator placement preserves it. & \stateactive \\

\commit{eac6e40}, \code{tx\_context.rs} & Cache-pad per-worker transaction atomics &
Workers update their own liveness/commit slots frequently. & \gainbounded: prevents multiple
workers' hot atomics sharing one cache line. & More memory per worker and a fixed worker
capacity. & \stateactive \\

\code{SnapshotCache} and thread-local worker identity & Cache snapshot/commit-log lookup
state per worker & Avoid repeated worker discovery and repeated full commit-log searches
during visibility checks. & \gainbounded: O(1) inline worker slots and cached lookup position;
no isolated benchmark. & \code{MAX\_WORKERS\_CAP=256}; stale cache state would return an
incorrect visibility answer. & \stateactive \\
\bottomrule
\end{longtable}

\section{Page layout and structural maintenance}

\begin{longtable}{IORGKS}
\toprule
\textbf{Commit / location} & \textbf{Optimization} & \textbf{Reason} &
\textbf{Gain / evidence} & \textbf{Potential risk or trade-off} & \textbf{State} \\
\midrule
\endhead
\commit{34555e2} & Delay internal-page overflow until one slot remains & The former
three-slot reserve triggered SMOs too eagerly. & \gainbounded: uses two additional key slots
before splitting, reducing premature structural work. & Too little headroom can leave no
space for the pending parent entry; boundary arithmetic must remain exact. & \stateactive \\

\commit{b29d347} & Tune \code{FAN\_OUT/NUM\_RECORDS} from 125 to 123 for a real 4 KiB
allocation & The wrapper around a block made the nominal size cross a page boundary. &
\gainmeasured: page-aligned allocations rose from 25\% to 100\%; dTLB load misses fell
9--12\%. & Two fewer records per base leaf can increase split frequency; result depends
on the exact key/payload layout. & \stateactive \\

\commit{1f6d666}, \code{BigTreeSize} & Give hot Warehouse/District tables independent
page-size classes & Tiny hot tables otherwise compact/version-split frequently; table
specific sizing avoids inflating every tree. & \gainmeasured: District root restarts reportedly
18.9M to 4.3M at the selected 32 KiB default. Short tpmC tests were broadly flat. & Large
leaves retain more dead versions and make aged scans slower; 512 KiB was measured about
4--6x slower than 32 KiB for delayed scans. & \stateactive \\

\commit{2a2607e} & Treat survivor count equal to capacity as requiring a key split & A
compacted page at exactly 100\% capacity cannot accept the pending write. & \gainenabling:
removes a root-leaf overflow panic and non-root infinite retry/livelock. & A key split may
copy more data than version compaction; off-by-one regressions are catastrophic. &
\stateactive \\

\commit{dc30111} & Tiered nearest-key-boundary search & A greedy nearest cut could miss a
farther boundary whose two halves both fit. & \gainenabling: guarantees a capacity-fitting
cut when one exists and avoids futile split retries. & Last-resort tiers may tear a
same-key historical chain or produce an imbalanced split when no ideal boundary exists. &
\stateactive \\

\commit{dcb5abe} & Merge with the emptier adjacent sibling & Always choosing the right
sibling ignored available space on the left. & \gainbounded: increases merge success and
reduces immediate re-splits. & Inspecting both candidates adds validation races; chosen
candidate may retire before upgrade. & \stateactive \\

\commit{4c01360} & Preserve append order when combining leaf records & Sorting by
transaction start time was unnecessary and could violate newest-physical-record lookup. &
\gainbounded: removes a sort and prevents wrong \code{rfind} results. & Correctness relies on
both inputs already preserving their append/version order. & \stateactive \\

\commit{4bc78de} & Count true GC survivors before merging & Active count omitted dead
predecessors still required for abort. & \gainenabling: avoids overflowing the fixed leaf;
falls back to key split when survivors do not fit. & Survivor evaluation must match the
later copy exactly or the pre-check is stale. & \stateactive \\

\commit{5c54658} & GC-gated survivor protection plus correct \code{bulk\_push} accounting &
Unresolved deletes need predecessors for abort; bulk copy incorrectly assumed all rows
were active. & \gainenabling: fixed two reproduced TPC-C crashes and prevents lost keys and
wrong capacity decisions. & Over-protection retains garbage; under-protection makes abort
silently lose data. & \stateactive \\

\commit{fe03d90} & Derive internal-entry liveness from disjoint live intervals & Remove a
separately mutated obsolete/version-region atomic from published child entries. &
\gainbounded: one atomic metadata field and mutation protocol removed. & Depends on live
sibling intervals remaining pairwise disjoint; overlapping publication breaks selection. &
\stateactive \\
\bottomrule
\end{longtable}

\section{Transactional version-chain optimizations}

\begin{longtable}{IORGKS}
\toprule
\textbf{Commit / location} & \textbf{Optimization} & \textbf{Reason} &
\textbf{Gain / evidence} & \textbf{Potential risk or trade-off} & \textbf{State} \\
\midrule
\endhead
\commit{bea1117}, update-in-place feature & Overwrite a committed payload when no live
snapshot can observe its predecessor & Avoid creating another MVCC version once GC proves
the old value unreachable. & \gainbounded: one tuple append and later compaction avoided per
eligible update. & Disabled with WAL because the physical shortcut is not represented by
one ordinary update record; an incorrect oldest-snapshot test violates SI. & \stateactive \\

\commit{85ac1b0}, \code{DbTransaction::update} & Self-overwrite a tuple already inserted
by this transaction & Repeated updates by one open transaction formerly appended one
unreclaimable tuple each. & \gainbounded: $N$ same-key updates create one replacement tuple,
so physical growth changes from O($N$) to O(1). & Must match the full \code{TxStamp}, not
only worker ID; WAL must still record every logical update needed for recovery. & \stateactive \\

Working-tree follow-up, \code{insert\_on\_tree} & Reuse the transaction-owned tuple across
repeated delete/reinsert cycles & After the first replacement, every later cycle was
another spelling of self-overwrite but still appended a tuple. & \gainbounded: $N$ cycles
create one replacement tuple; commit exposes the final payload, abort restores the
original, and WAL recovery reconstructs the final value. & Live/dead counters must be
reversed exactly; no extra write-set entry may be added; reuse is legal only when both
insertion and deletion stamps equal the current transaction. & \stateactive \\

\commit{160186e} & Abort repeated-key writes in strict LIFO order & Reverting only the
newest tuple left older self-written state visible. & \gainenabling: complete rollback across
all same-key and cross-table writes; also makes tuple-reuse write-set semantics safe. &
Readers may observe per-key rollback progress during abort; changing order can expose an
intermediate version. & \stateactive \\

First-writer-wins conflict detection & Stop another transaction extending an unresolved
same-key chain & Bound concurrent pending versions and retain OSIC's conflict rule. &
\gainbounded: at most the owning transaction can mutate the unresolved live key. & Aborts
under contention waste traversals/work; TPC-C conflict rates remain a major performance
opportunity. & \stateactive \\
\bottomrule
\end{longtable}

\section{Garbage collection and reclamation}

\begin{longtable}{IORGKS}
\toprule
\textbf{Commit / location} & \textbf{Optimization} & \textbf{Reason} &
\textbf{Gain / evidence} & \textbf{Potential risk or trade-off} & \textbf{State} \\
\midrule
\endhead
\commit{2e3ad21} & Non-blocking dead-page reuse checks & Blocking allocation locks could
join the page-latch dependency graph and deadlock. & \gainenabling: allocator can bail out and
allocate fresh memory rather than block a structural operation. & Falling back increases
memory use when reuse metadata is contended. & \stateactive, evolved \\

\commit{95e149f} then current \code{TxContext} & Replace global active-transaction mutex;
ultimately use per-worker liveness slots & Global minimum-snapshot bookkeeping serialized
transaction entry/exit. & \gainbounded: liveness publication is per-worker atomic; minimum is
computed without one global transaction lock. & Scanning fixed slots costs O(workers);
registration ordering must prevent a just-starting reader being missed. & \stateactive \\

\code{block\_tracer.rs}, worker-sharded \code{SkipMap} & Shard retired-block queues & A
single dead-page list/index was a reclamation hotspot. & \gainbounded: independent inserts;
allocators check their own shard first and steal round-robin only on demand. & Uneven
shards can retain reusable memory; stealing adds cross-shard traversal. & \stateactive \\

\commit{f9fc960} & Own-shard-first reuse with round-robin stealing & Improve cache
locality and avoid eagerly touching every reclamation shard. & \gainbounded: common reuse
touches one shard; cross-shard work occurs only on a miss. & Local starvation or imbalance
can delay reuse until stealing runs. & \stateactive \\

\commit{d0e0977} & Fold retirement into SmartCell's version word & Separate retired state
raced guard Drop and block reuse. & \gainenabling: eliminated a reproduced GC hang/livelock
and one atomic field/state transition. & Bit packing makes latch arithmetic delicate;
unlock must preserve retirement bits. & \stateactive \\

\commit{ecc9900} & Key retired pages by (death version, address) & Two siblings retired by
one merge share a death version and collided in the old index. & \gainenabling: prevents one
of the two blocks leaking on every such collision. & Address participates only in
identity; reuse must not compare it as a temporal value. & \stateactive \\

\commit{ecc9900} & Atomic insertion/deletion stamps & Lock-free readers could tear a
plain stamp while a writer deleted or invalidated it. & \gainenabling: ThreadSanitizer race
removed without adding reader locks. & Relaxed ordering solves atomicity, not publication;
payload publication still relies on the node latch. & \stateactive \\

Historical-block routing plus \code{live\_min\_snapshot} reuse gate & Keep old readers on
retired source blocks instead of copying all old-reader-visible records forward & Copying
for the oldest reader globally couples all keys and can self-reinforce into compaction
livelock. & \gainbounded: replacement pages retain only structurally necessary survivors;
old blocks become reusable immediately after their last reader. & Reusing a block before
the oldest routed reader exits is a use-after-reuse; root/version routing must select the
correct historical block. & \stateactive \\
\bottomrule
\end{longtable}

\section{Root indexing and database sharding}

\begin{longtable}{IORGKS}
\toprule
\textbf{Commit / location} & \textbf{Optimization} & \textbf{Reason} &
\textbf{Gain / evidence} & \textbf{Potential risk or trade-off} & \textbf{State} \\
\midrule
\endhead
\commit{816356a}, \code{FrugalList} & Compact append-oriented root-version index & Root
lookups need predecessor-by-version while new structural roots append over time. &
\gainbounded: compact linear representation with cheap append, suited to few roots. & Lookup
becomes linear as history grows; not ideal for long GC-off runs. & \stateactive option \\

\commit{24fed9b} and later wrappers & Runtime-selectable Frugal, linked-list, skip-list,
and B-tree root indexes & Permit experiments across append cost, memory, and predecessor
lookup complexity. & \gainenabling: chooses an index appropriate to root-history size; no
single controlled winner retained. & Multiple implementations increase validation and
maintenance surface; semantics must match at version boundaries. & \stateactive options \\

\commit{5f92e9f} & One tree/index per database table & Encoding table IDs into one shared
tree creates unrelated contention and oversized structural domains. & \gainbounded: table
writes and SMOs proceed independently; table-specific capacities become possible. &
Cross-table atomicity moves to transaction/WAL coordination; more roots and allocators. &
\stateactive \\

\commit{0ba6ca5}/\commit{73cc4a5} & Track live transactions without GC and optionally
truncate commit logs & Visibility needs liveness even when block GC is disabled; unbounded
commit history consumes memory. & \gainbounded: truncation bounds retained commit-log history
when configuration permits. & Pruning before every relevant reader has advanced causes
false invisibility; long historical readers delay reclamation. & \stateactive \\
\bottomrule
\end{longtable}

\section{Range scans and analytical terminals}

\begin{longtable}{IORGKS}
\toprule
\textbf{Commit / location} & \textbf{Optimization} & \textbf{Reason} &
\textbf{Gain / evidence} & \textbf{Potential risk or trade-off} & \textbf{State} \\
\midrule
\endhead
\commit{816356a}/\commit{f4e786b} and current iterator & Lazy range iteration & Avoid
materializing the entire range result before consumption. & \gainbounded: memory is one path
and at most one leaf buffer instead of O(result size). & Iterator must revalidate across
SMOs and must release its registered snapshot on every exit path. & \stateactive \\

Current cursor-based internal routing & Select the newest visible child containing the
scan cursor, without \code{unique\_by} or child lists & Physical internal slots contain
overlapping historical routes; equality deduplication was both costly and incorrect. &
\gainbounded: no child vector, hash set, sort, or two \code{unique\_by} adaptors; GC stress no
longer returns duplicate keys. & Cursor increment saturates at maximum key and needs an
explicit completion test; wrong newest-route ordering causes duplicates or omissions. &
\stateactive \\

\code{for\_each\_ref}/\code{fold\_ref} streaming terminals & Pass callback-scoped payload
references instead of owned result rows & Analytical consumers immediately aggregate and
do not need cloned handles or result objects. & \gainmeasured: full scan 1.165--1.173x faster,
about 16.8\%; 28.8--29.3M versus 24.7--25.1M rows/s. & Callback cannot retain the payload
or safely re-enter the same tree; long callbacks prolong leaf access. & \stateactive \\

\code{try\_for\_each\_ref} & Fallible early-exit visitor & Stop scans when a consumer has
enough data or encounters an error. & \gainbounded: avoids visiting all remaining records;
gain depends on stop position. & Drop must release iterator-owned snapshots but not a
transaction-owned snapshot. & \stateactive \\

\code{count\_ref}/\code{range\_count} & Specialized visible-row count & Counting does not
need \code{RecordPointResult} construction or payload-handle cloning. & \gainmeasured:
approximately 49--53M rows/s in the recorded full-scan count benchmark; retained after
unsafe routing experiments were reverted. & Must use identical routing and visibility
semantics to the normal iterator. & \stateactive \\

\commit{8e48ae5} & Test key range before MVCC visibility & Narrow scans can cheaply reject
out-of-range physical rows before up to two visibility lookups. & \gainmeasured: CH Q4 about
18\% faster; Q5 about 11--13\% faster in short indicative runs; no expected full-scan
gain. & Benefit depends on selectivity; changing predicate order is safe only because both
checks are side-effect-free for accepted rows. & \stateactive \\

\commit{8e48ae5} & Monomorphize the range iterator visibility closure & Remove a
\code{dyn FnMut} indirect call for every tested insertion/deletion stamp. & \gainbounded:
indirect call becomes inlinable; incremental gain was below run-to-run noise after filter
reordering. & Exposing raw snapshot cache/log slices broadens a sensitive internal API;
visibility logic must remain identical. & \stateactive \\

Working tree, \code{iter\_query.rs} & Exact full-domain leaf-loop specialization & Q1/Q6
do not need to compare every physical key with minimum and maximum bounds. & \gainbounded:
removes two bound comparisons per physical record; included in the combined YCSB-E result
below, without isolated attribution. & Only the exact \code{Key::MIN..=Key::MAX} interval
may take this path; bounded ranges need their range predicate. & \stateactive \\

\commit{c0fe172} & Search physical leaf order rather than assuming transaction-start
ordering & Append order is not globally sorted by \code{ts\_start}; binary/search-like
selection could miss the newest matching row. & \gainenabling: fixes point and range misses;
also avoids an invalid ordering assumption. & Linear reverse search costs O(page size),
making compaction/page sizing important for scan performance. & \stateactive \\
\bottomrule
\end{longtable}

\section{WAL, durability, and encoding}

\begin{longtable}{IORGKS}
\toprule
\textbf{Commit / location} & \textbf{Optimization} & \textbf{Reason} &
\textbf{Gain / evidence} & \textbf{Potential risk or trade-off} & \textbf{State} \\
\midrule
\endhead
\commit{3e8b837}/\commit{8a39d51} & Group-commit linger & Allow the background writer to
collect more records per hardening operation. & \gainbounded: amortizes write/fsync across a
larger batch; numeric isolated gain not retained. & Adds commit latency and requires
correct wait/wakeup handling. & \stateactive \\

\commit{ab10660} & Non-blocking WAL enqueue with background batch hardening & Committers
should not serialize on file I/O. & \gainbounded: caller pays enqueue rather than pwrite/fsync;
durability advances asynchronously. & Backpressure and hardened-watermark correctness;
unbounded queues may grow if storage falls behind. & \stateactive \\

\commit{4a1a50a} & Table-driven CRC32 & Bit-at-a-time framing checksum consumed 22.6\% of
YCSB-A CPU cycles. & Combined three-fix campaign: CRC self-time fell 22.6\% to 11.3\%; it
shares the end-to-end numbers below with build and sharding changes. & Must preserve the
same polynomial and on-disk format; table lookup can affect cache footprint. & \stateactive \\

\commit{4a1a50a} & Four sharded WAL producer channels and flush threads & One crossbeam
channel spent about 8\% of total TPC-C CPU in producer spin-wait. & Combined campaign:
TPC-C +7.4\%/+17.6\% at 16/64 threads; YCSB-A +25.4\%/+68.7\%; YCSB-E +16.2\%/+49.8\%;
YCSB-F +20.7\%/+70.1\%. & Multiple flushers complete out of order; a single max watermark
was unsound and required per-shard submitted/confirmed tracking. & \stateactive \\

\commit{4a1a50a} & Make \code{target-cpu=native} effective in benchmark builds & The old
configuration silently failed to apply intended CPU-specific optimization. & Part of the
same combined campaign; no independent percentage. & Native binaries are not portable to
older/different CPUs and benchmark comparisons must use identical build flags. & \stateactive \\

\commit{ebdd789}, \code{LockFreeWalWriter} & Reserve file offsets atomically and let each
worker \code{pwrite} its own local batch & Remove the shared producer channel/background
writer as an alternative backend. & \gainmeasured synthetic at 16 threads: batch-64
5.09M ops/s versus sharded writer 1.35M; unbatched only 0.75M and loses. & Shared-file
extension and syscalls dominate without batching; recovery and watermarks must handle
out-of-order writes. & \stateactive option \\

\commit{9ced1a3} & Per-shard submitted/confirmed durability aggregation & Raw max let a
fast shard falsely cover an unflushed slow shard; raw min could wait forever on an idle
shard. & \gainenabling: fixed recovery returning 1200 rather than 1800 records and eliminated
the alternative hang; failing test passed 5 times and full suite 3 times. & Submission
must publish before enqueue; stale underestimates allow premature durability. & \stateactive \\

\commit{efc3773} & Ticket-less production WAL logging & Every write allocated a bounded
ack channel whose receiver production callers discarded. & Combined with size hints:
TPC-C +12.6\%, YCSB-A +25.3\%, synthetic 16-thread writer +79\%. & Tests/callers that
really wait for one record must use explicit \code{*\_with\_ticket} methods. & \stateactive \\

\commit{efc3773} & Payload-aware exact framing capacity hints & Buffers sized for 8-byte
tests repeatedly grew for 300--1000-byte real rows. & Same combined result; allocator
self-time \code{RawVecInner::finish\_grow} fell from 7.1\% to 3.0\% in TPC-C. & Every
payload encoder must keep its hint consistent; underestimates lose performance, while
overestimates waste memory. & \stateactive \\

\commit{73cc4a5}/\commit{a5e6b7a} & Commit-log truncation configuration & Long runs with
GC disabled otherwise retain every commit forever. & \gainbounded: bounds memory once all
relevant snapshots pass a segment. & Premature truncation makes committed versions look
uncommitted; long readers pin history. & \stateactive \\
\bottomrule
\end{longtable}

\section{Workload and benchmark-path optimizations}

\begin{longtable}{IORGKS}
\toprule
\textbf{Commit / location} & \textbf{Optimization} & \textbf{Reason} &
\textbf{Gain / evidence} & \textbf{Potential risk or trade-off} & \textbf{State} \\
\midrule
\endhead
\commit{423a1f5} & Fast YCSB RNG, reference-counted rows, and pre-sized scan buffers &
Benchmark-driver allocation/RNG overhead obscured engine throughput. & \gainbounded: fewer
row clones, reallocations, and expensive random-number operations; no isolated percentage
retained. & Driver optimizations can change the generated distribution or accidentally
benchmark shared payload handles rather than realistic copies. & \stateactive \\

\commit{eac6e40} & Use \code{fastrand} for TPC-C data and transaction parameters & Random
generation was on the transaction hot path. & \gainbounded: cheaper parameter generation; no
isolated percentage. & Must retain TPC-C distributions and invalid-item rollback rate. &
\stateactive \\

\commit{47ec877} & Scramble Zipfian rank with FNV & Sequential rank mapping created a
key-locality pattern unlike reference YCSB/LeanStore. & \gainenabling: more representative
contention and comparison, not claimed as an engine speedup. & Changes benchmark results;
before/after engine numbers are incomparable unless generator version is fixed. & \stateactive \\

Working tree, \code{ycsb\_random.rs} & Walker alias Zipf sampler & Remove floating-point
power from every timed request. & \gainmeasured: part of the combined YCSB-C improvement
from 7.16M to 14.82M ops/s (+106.8\%). & O(n) setup and 8 bytes/key (about 160 MB at 20M
keys); \code{f32} probabilities slightly quantize the distribution. & \stateactive \\

Working tree, \code{ycsb\_schema.rs}/\code{ycsb\_txn.rs} & In-place bulk value generation
and \code{writeallfields} policy & Avoid temporary value containers and implement standard
YCSB single-field updates. & \gainmeasured: full-row YCSB-A 2.52M to 3.11M (+23.5\%);
standard single-field mode reached 4.81M ops/s (+90.8\% versus the old full-row baseline).
& The larger number changes semantics; historical comparison requires
\code{write\_all\_fields=true}. Atomic patch closure runs under a leaf write latch. &
\stateactive \\

Working tree, \code{query.rs} & \code{point\_exists\_si} & Hit/miss-only YCSB reads do not
need result vectors or payload-handle clones. & \gainmeasured: included with alias sampling
in YCSB-C's 7.16M to 14.82M ops/s combined result. & Cannot serve callers that need payload
contents; reader registration and SI visibility must remain gap-free. & \stateactive \\

Working tree, transaction commit paths & Elide empty read-only commits & Avoid advancing
the global clock and appending a commit-log entry when no table was written. &
\gainbounded: one timestamp increment and commit-log operation removed per empty read-only
transaction. & Successful read-only commits return no commit timestamp; callers must honor
the existing \code{Option<Version>} contract. & \stateactive \\

Working tree, \code{ycsb\_driver.rs} & Sample scan latency and amortize timeseries clocks &
Measure one of every 1,024 scans and refresh the elapsed-time bucket every 256 operations.
& \gainmeasured: combined YCSB-E throughput 4.59M to 6.22M ops/s (+35.4\%). & CSV latency
count is sample count; percentiles have sampling error and bucket placement can lag by 255
operations per worker. & \stateactive \\

Working tree, \code{olap\_scan.rs}/engine wrapper & Isolated CH Q1 and Q6 modes & Prevent
Q1/Q6 benchmark points from paying for the unrelated Q4/Q5 joins in the old rotating mode.
& \gainenabling: results now attribute work to the named query; no isolated speedup claimed.
& New isolated results measure a different workload and are not comparable to historical
four-query \code{ch} runs. & \stateactive \\

Working-tree benchmark scaling & Run TPC-C/YCSB/WAL comparisons as ordinary smoke tests;
restore full scale with \code{CMVBT\_FULL\_BENCH=1} & Previously ignored/full benchmarks
did not run on normal developer machines. & \gainenabling: complete suite now runs 138 tests,
zero ignored, in 8.23 seconds on the development machine. & Smoke scale catches hangs and
gross regressions, not production throughput or long-duration memory growth. & \stateactive \\

\code{RESTART\_TRACE=false} plus reset per run & Prevent diagnostic strings accumulating
across repeated benchmark invocations & A checked-in tracing flag grew past 12 GiB and
OOM-killed the WAL comparison. & \gainenabling: removes unbounded process-lifetime trace
growth; tiny repro stayed bounded after reset. & Turning tracing on still has high
overhead; reset must run before every repeated workload. & \stateactive \\
\bottomrule
\end{longtable}

\section{Rejected, reverted, or deliberately optional experiments}

These entries are part of the optimization history but are not unconditional current
wins. Keeping them visible prevents rediscovering an unsafe shortcut without its context.

\begin{longtable}{IORGKS}
\toprule
\textbf{Commit / location} & \textbf{Attempt} & \textbf{Why attempted} &
\textbf{Observed result} & \textbf{Failure mode / risk} & \textbf{State} \\
\midrule
\endhead
\commit{d565b2d}--\commit{5c26abb} & ORWC optimistic read/write coupling & Seek
non-blocking upgradeable writers alongside free readers. & No trustworthy gain retained;
about 460 lines were ultimately removed. & Repeated upgrade races, dangling readers, and
eventual blocking made it duplicate OLC without a sound distinct guarantee. & \statereverted \\

\commit{45cb048} & Hold one long-lived range-scan block guard & Avoid re-borrowing blocks
during iteration. & Not retained. & GC did not know about the guard lifetime; reuse could
invalidate the long-lived reference. & \statereverted \\

Range items 1/2/4/5 in \code{range\_scan\_iteration.md} & Eager child-list reuse, small-page
linear dedup, hybrid traversal, and lean frames & Remove allocation/hash work and release
internal pages early. & Experimental single-set version appeared 18--20\% faster. & Equal
boundaries cannot prove arbitrary historical intervals superseded; GC stress returned
duplicate logical keys. & \statereverted \\

\commit{0659d34} experiment & Retry alternate merge candidate immediately & Increase
merge success when the first sibling cannot be upgraded. & No gain retained. & Tight
retry loop lacked backoff and starved workers under contention. & \statereverted \\

Overflow-headroom companion to \commit{85ac1b0} & Split leaves earlier to reserve more
compaction space & Reduce the chance a pending write meets a completely full leaf. & No
valid gain; stress hung. & More SMOs amplified contention while the actual off-by-one bug
remained. & \statereverted \\

Per-reader LCB survivor checks & Copy a predecessor if any active reader might need it &
Attempt to avoid one global oldest-reader watermark. & Repro fired 517,748 checks in 12
seconds for one fixed reader; no progress. & An open transaction may read any key, so the
predicate protects nearly everything and creates self-reinforcing compaction livelock. &
\statereverted \\

Thread-local raw lock-free WAL buffer & Reuse encoding memory on the caller thread & Avoid
one allocation per write. & No production gain: real calls already enter persistent
\code{LocalBatch} buffers. & Dead code/complexity benefiting only the synthetic unbatched
test path. & \statereverted \\

\commit{6486d97}, mimalloc feature & Make cross-thread frees cheaper for sharded WAL &
WalWriter allocates on producers and frees on flush threads. & YCSB-A +3.0\% at 8 threads,
+2.3\% at 16; lock-free backend within noise. TPC-C fell 3.3--4.9\%. & Workload-dependent
allocator trade-off; jemalloc remains the default. & \stateoptional \\

Very large \code{BigTreeSize} leaves & Defer compaction on the hottest tiny tables &
Short tpmC results were flat and restart counts improve with size. & Aged 512 KiB scans
were 4--6x slower than 32 KiB because they touch more dead records. & \stateoptional \\

Unbatched \code{LockFreeWalWriter} & Remove all channels and batching & Simplest direct
atomic-offset append. & 0.75M ops/s at 16 threads versus 1.35M for sharded WalWriter and
5.09M for local batch-64. & One syscall per logical WAL record and shared inode-extension
serialization. & \stateoptional, not recommended \\
\bottomrule
\end{longtable}

\section{Combined measured gains at a glance}

This final table repeats only repository-backed quantitative results. It is intentionally
short and is the correct table to cite when a numerical gain is required.

\begin{longtable}{p{6.1cm}p{5.2cm}p{5.5cm}p{6.6cm}}
\toprule
\textbf{Optimization set} & \textbf{Workload} & \textbf{Measured change} &
\textbf{Qualification} \\
\midrule
\endhead
4 KiB base-block tuning & Allocation / TLB profile & Alignment 25\% to 100\%; dTLB load
misses $-9$ to $-12$\% & Exact object layout dependent. \\
CRC32 + effective native build + four WAL shards & TPC-C, 16/64 threads & +7.4\% / +17.6\%
& Combined result; do not assign all gain to one component. \\
Same three WAL/build changes & YCSB-A, 16/64 & +25.4\% / +68.7\% & Short, single campaign;
write-heavy Zipfian workload. \\
Same three WAL/build changes & YCSB-E, 16/64 & +16.2\% / +49.8\% & Scan-heavy but with a
write component. \\
Same three WAL/build changes & YCSB-F, 16/64 & +20.7\% / +70.1\% & Read-modify-write. \\
Current YCSB/scan batch & YCSB-A full-row / standard single-field & 2.52M to 3.11M
(+23.5\%) / 4.81M (+90.8\% vs old baseline) & 200k rows, 4 threads, 3 seconds; second
number uses standard, different update semantics. \\
Alias Zipf + hit-only point path & YCSB-C & 7.16M to 14.82M ops/s (+106.8\%) & Combined
batch result on the same short 200k-row setup. \\
Scan-loop and measurement batch & YCSB-E & 4.59M to 6.22M ops/s (+35.4\%) & Combined
batch result; latency samples one scan in 1,024. \\
Ticket-less WAL + exact buffer hints & TPC-C / YCSB-A & +12.6\% / +25.3\% & Same short
profiling setup before and after both fixes together. \\
Ticket-less WAL + exact buffer hints & Synthetic WalWriter, 16 threads & +79\% & 8-byte
payload makes removed ticket allocation unusually dominant. \\
Lock-free local batch-64 & Synthetic WAL, 16 threads & 5.09M vs 1.35M ops/s & Debug-build
microbenchmark; unbatched lock-free is slower. \\
Zero-copy streaming range terminal & 200k-row full scan & 1.165--1.173x; about +16.8\% &
40 alternating samples after warmup. \\
Range-before-visibility predicate & CH Q4 / Q5 & about +18\% / +11--13\% & Indicative
short contended runs; benefit is for narrow ranges only. \\
Specialized count terminal & Full scan & about 49--53M rows/s & Absolute throughput, not a
recorded controlled percentage versus current iterator. \\
BigTreeSize 32 KiB default & District root restarts & 18.9M to 4.3M & Commit-reported,
not reproduced in the later size sweep. \\
Mimalloc feature & YCSB-A / TPC-C & +2.3--3.0\% / $-3.3$ to $-4.9$\% & Kept optional;
not a universal optimization. \\
Scaled complete test/benchmark smoke suite & Development test run & 138 passed, 0 ignored,
8.30 s & Correctness/operability result, not production performance. \\
\bottomrule
\end{longtable}

\section*{Maintenance rule}

When adding an optimization, add one row to the appropriate detailed table and, only if
a controlled number exists, one row to the measured-gains table. Record the workload,
thread count, baseline, resulting value, and whether multiple changes were measured
together. For correctness/bounded-growth changes, state the exact eliminated cost instead
of manufacturing a percentage. Reverted experiments remain in the rejected table with
their failure mode.

\end{document}
